\begin{mvtable}
\phantomsection\label{tab:5.7}\mvtabcap{表 5.7\quad 音乐才能测试数据的 $95\%$ 单一区间、庞弗罗尼区间及 $T^{2}$ 区间}
\begin{center}\small
单一置信区间用 $t_{95}(0.025)=1.99$\\
庞弗罗尼区间用 $t_{95}(0.025/7)=2.75$\\
联合 $T^{2}$ 区间(即置信椭球的投影)用 $F_{7,95}(0.05)=2.11$
\end{center}
\begin{tabular}{lcccccc}
\toprule
\multirow{2}{*}{变量} & \multicolumn{2}{c}{单一区间} & \multicolumn{2}{c}{庞弗罗尼区间} & \multicolumn{2}{c}{椭球的投影}\\
\cmidrule(lr){2-3}\cmidrule(lr){4-5}\cmidrule(lr){6-7}
 & 下限 & 上限 & 下限 & 上限 & 下限 & 上限\\
\midrule
$X_1$=旋律 & 26.93 & 29.27 & 26.48 & 29.72 & 25.76 & 30.44\\
$X_2$=和声 & 25.41 & 27.79 & 24.96 & 28.24 & 24.23 & 28.97\\
$X_3$=速度 & 34.63 & 36.17 & 34.33 & 36.47 & 33.85 & 36.95\\
$X_4$=拍子 & 33.16 & 35.24 & 32.76 & 35.64 & 32.12 & 36.28\\
$X_5$=短句 & 22.84 & 24.36 & 22.54 & 24.66 & 22.07 & 25.13\\
$X_6$=协调 & 21.20 & 22.80 & 20.90 & 23.10 & 20.41 & 23.59\\
$X_7$=风格 & 21.88 & 23.52 & 21.57 & 23.83 & 21.07 & 24.33\\
\bottomrule
\end{tabular}
\end{mvtable}

比较\hyperref[tab:5.7]{表 5.7} 与\hyperref[tab:5.6]{表 5.6},可以发现,\hyperref[tab:5.7]{表 5.7} 中的区间要更大一些. 不过对较大的样本容量 $n=96$,置信限仅在数据的第三位数字(即十分位数)上有典型的差别.

\section{多元质量控制图}
\phantomsection\label{sec:5.6}

为改进产品和服务的质量,需要考虑数据以查明引起质量变动的原因. 当一个制造过程为连续地生产产品,或当我们监控某种服务活动时,就应收集数据,用以评估这个过程的能力与稳定性. 当一个过程稳定时,变动是由一些持续存在的常见原因引起的,没有一个原因是变动的主要来源.

任何控制图的目的,就是为了识别是否出现了引起变动的特殊原因,这些原因来自常规过程的外部. 这些变动原因的出现表明需要某种适时的维修,不过也可能暗示需要改进此过程. 控制图可使这些变动可视化,因而有利于我们区分引起变动的一般原因和特殊原因.

控制图一般由按时间顺序标绘的数据点和两条水平线组成,这两条水平线称为控制限,用来表示由一般原因引起的变动量. 有一种有用的控制图称为 $\bar{X}$ 控制图,其生成过程如下:

1. 按时间顺序对各个观测值或样本均值作出标绘.

2. 生成并画出代表所有观测值的样本均值的中心线 $\bar{x}$.

3. 按如下公式计算并画出控制限:

\begin{center}
控制上限 (UCL) $=\bar{x}+3\times$ 标准差\\
控制下限 (LCL) $=\bar{x}-3\times$ 标准差
\end{center}

控制限中的标准差可用被标绘观测值的标准差来估计. 对单一的观测值,上述标准差往往就是样本标准差. 假如所标绘的是容量为 $m$ 的若干个样本的均值,则其标准差为样本标准差除以 $\sqrt{m}$. 控制限中之所以选取加减 3 倍标准差,是为了在数据服从正态分布的假设下,使误将正常数据标成失控数据的可能性变得非常小. 所谓失控数据,是指暗示存在某种特殊变动原因的观察值.

\noindent\textbf{例 5.8}\quad (构造一元控制图)

威斯康星州麦迪逊警察局把对自身行为的定期监测作为其管理质量改进计划的一部分. \hyperref[tab:5.8]{表 5.8} 给出了 5 类加班小时数的数据,每一个观察值代表 12 个工资支付期(大约为半年)内加班小时数之和.
